% Cover letter for the transfer to Applied Stochastic Models in Business and Industry (Wiley).
\documentclass[11pt]{article}
\usepackage[utf8]{inputenc}
\usepackage{cmap}
\input{glyphtounicode}
\pdfgentounicode=1
\usepackage[margin=1in]{geometry}
\usepackage{mathptmx}
\usepackage[hidelinks]{hyperref}
\usepackage{xurl}
\usepackage{parskip}

\newcommand{\venue}{Applied Stochastic Models in Business and Industry}

\begin{document}

\thispagestyle{empty}

\noindent
The Editors\\
\textit{\venue}

\bigskip
\noindent \today

\bigskip
\noindent Dear Editors,

\noindent I submit the manuscript ``Online Anchor-Robust Forecasting under Economic
Regime Shift'' for consideration as a research article in \textit{\venue}. The paper
comes to you through the Wiley Transfer Desk. It was reviewed at the
\textit{Journal of Forecasting} (manuscript 7441141) and declined there. I have
since revised it substantially, and the version uploaded here replaces the one that
was reviewed. A separate document lists each point of the reviewer's report and the
change made in response.

The paper addresses a practical forecasting problem: a model fitted on past regimes
can fail when a driver such as geopolitical risk or policy uncertainty moves to a
level not seen before. Most online methods react after the loss has risen. The paper
asks whether a forecaster can be made robust in advance when the shift moves an
observed driver. The method, OARF, is an online version of anchor regression with a
single penalty that reduces it to online gradient descent when set to zero. For a
supplied driver channel I prove a no-regret bound against a comparator evaluated
under the worst bounded shift of the driver. A simulation study with a known
structural model shows a large gain in worst-case error and uses stress tests to
show where the gain disappears. An application to weekly Brent crude and Henry Hub
gas volatility shows that the learned combinations beat a persistence forecast but
that OARF is not better than online gradient descent there, and the paper reports
this result as it is.

I think the paper fits the journal because it combines a stochastic model of regime
change, an online estimator with a guarantee, and a real business application, and
because it gives a practitioner a concrete procedure for deciding when the method
should and should not be used.

The main changes since the review are the following. The real-data target is now
weekly realised variance over non-overlapping weeks, so no part of the target is
known when the forecast is made; the earlier overlapping daily design is kept in an
appendix for comparison. A persistence forecast and an equal-weight average are
reported as benchmarks. The paper now opens with a short worked forecasting example,
presents the basic method with a known anchor before any extension, states what the
anchor must satisfy in the empirical application, and moves the learned-anchor
extension to its own section after the main results.

Code, data-build scripts and the commands that regenerate every number, table and
figure are available at
\url{https://github.com/vinhqdang/OARF-Online-Anchor-Robust-Forecasting-under-Economic-Regime-Shift}.

The manuscript is original, has not been published, and is not under
consideration elsewhere. I have no competing interests to declare. The study uses
only public, non-personal data and did not require ethics approval.

\noindent Thank you for considering the manuscript.

\bigskip
\noindent Sincerely,\\[2pt]
Quang-Vinh Dang\\
British University Vietnam, Hung Yen, Vietnam\\
\href{mailto:vinh.dq4@buv.edu.vn}{vinh.dq4@buv.edu.vn}

\end{document}
